% =====================================================================
% Section 2: related work (about 1.25 columns of text by a glyph-width
% estimate, plus Table 1; the outline target was ~0.9 column).
% Table 1 (tab:positioning) lives in tables/tab_positioning.tex and is
% included at the top so that it floats next to the section. Each
% paragraph ends with what the line of work lacks for CST and what we
% reuse. Sec. 1 already outlines the SimulST lineage and the four closest
% systems (JSTAR, DiariST, STAC-ST, Seed LiveInterpret 2.0); this section
% adds the details and Table 1 lists the differences. Latency metrics are
% defined in Sec. 4, corpora are compared in App. A (tab:corpora), and
% baselines are run in Sec. 6.
% If space is short, cut first: the transducer clause of the last
% paragraph, then the code-switching sentence, then "and Hibiki's
% contextual alignment". Already cut for space (re-add if room): a
% MuST-C/IWSLT benchmark sentence (Sec. 1 makes that point) and
% "Diarization benchmarks use monolingual calls or meetings" with
% cieri2004fisher, canavan1997callhome and vinnikov2024notsofar.
% =====================================================================
\section{Related Work}
\label{sec:related}

% =====================================================================
% Table 1 (tab:positioning): CST against the closest settings and systems.
% Single column, \scriptsize, 10 columns; estimated natural width about
% 228pt with \tabcolsep=2pt (ICML column: 234.87pt). If it overflows,
% lower \tabcolsep to 1.5pt before anything else.
% Included from Sec. 2 (sec:related); every row is cited in Sec. 1 or 2.
% Cells checked against the cited papers: SeamlessStreaming is evaluated
% on FLEURS; Hibiki and Hibiki-Zero release weights; STAC-ST's merged
% Fisher-CALLHOME test channels contain cross-talk (11-12% overlap) and
% it releases scripts, not models; DiariST is scored on single-channel
% audio (first array channel, headset mix) with natural overlap and
% releases its evaluation set, offline baselines and scorer, not its
% streaming model; JSTAR uses a 5-channel smart-glasses array, tells
% wearer from partner by the direction of speech, tests on in-house real
% conversations and on FLEURS with simulated overlap, and releases
% nothing; Seed LiveInterpret 2.0 depicts a two-speaker Zh-En conversation
% (its Fig. 2) but is evaluated per direction on RealSI (annotations and
% video links public), does not report its input channels and is closed.
% \pmark = partial (some systems, or shown but not evaluated).
% =====================================================================
\begin{table}[t]
  \caption{\task against the closest settings and systems. Spk$\times$Lang: speakers $\times$ languages in one input stream; Auto dir.: the system chooses the direction of each turn; Ovl.: overlapping speech evaluated (sim.: simulated overlap); Str.: streaming; Out.: A~transcript, T~translation, D~speaker labels or turns, S~translated speech; Pub.\ data, Open wts: public evaluation data, open weights. Segm.: pre-segmented speech (e.g., MuST-C); long: unsegmented long-form audio (e.g., IWSLT 2025--2026). \cmark~yes, \xmark~no, \pmark~partial, n.r.~not reported. $^a$TAXI under its free scientific licence from BAS, rebuilt with our scripts and checksums; \syn under CC licences.}
  \label{tab:positioning}
  \centering
  \scriptsize
  \setlength{\tabcolsep}{2pt}
  \begin{tabular}{@{}lccccccccc@{}}
    \toprule
    Work & Input & \makecell{Spk$\times$\\Lang} & Dir. & \makecell{Auto\\dir.} & Ovl. & Str. & Out. & \makecell{Pub.\\data} & \makecell{Open\\wts} \\
    \midrule
    SimulST (segm.)   & mono   & 1$\times$1 & X$\to$Y               & \xmark & \xmark & \cmark & T     & \cmark & \pmark \\
    StreamST (long)   & mono   & 1$\times$1 & X$\to$Y               & \xmark & \xmark & \cmark & T     & \cmark & \pmark \\
    SeamlessStreaming & mono   & 1$\times$1 & X$\to$Y               & \xmark & \xmark & \cmark & T,S   & \cmark & \cmark \\
    StreamSpeech      & mono   & 1$\times$1 & X$\to$En              & \xmark & \xmark & \cmark & A,T,S & \cmark & \cmark \\
    Hibiki(-Zero)     & mono   & 1$\times$1 & X$\to$En              & \xmark & \xmark & \cmark & T,S   & \cmark & \cmark \\
    STAC-ST           & mono   & 2$\times$1 & Es$\to$En             & \xmark & \cmark & \xmark & A,T,D & \cmark & \xmark \\
    DiariST           & mono   & N$\times$1 & Zh$\to$En             & \xmark & \cmark & \cmark & T,D   & \cmark & \xmark \\
    JSTAR             & array  & 2$\times$2 & Es$\leftrightarrow$En & \cmark & sim.   & \cmark & A,T,D & \xmark & \xmark \\
    \makecell[l]{Seed\\LiveInterpret 2.0}
                      & n.r.   & \pmark     & Zh$\leftrightarrow$En & \pmark & \xmark & \cmark & T,S   & \cmark & \xmark \\
    \midrule
    \makecell[l]{\bench{} +\\\model{} (ours)}
                      & mono   & 2$\times$2 & \makecell{En$\leftrightarrow$De\\Ko$\leftrightarrow$En}
                                                                    & \cmark & sim.   & \cmark & A,T,D & \cmark$^a$ & \cmark \\
    \bottomrule
  \end{tabular}
  \ifshowdraftnotes
    \par\vspace{2pt}
    \parbox{\linewidth}{\raggedright\draftnote{planned in our row: the L1 overlap renders, \syn{} (incl.\ \pair{Ko}{En}), transcript and diarization scoring, and \model{} with its weights; the TAXI build scripts and the translation scorer exist.}}
  \fi
\end{table}


\paragraph{Simultaneous and streaming speech translation.}
Beyond fixed wait-$k$ \citep{ma2019stacl}, learned policies use monotonic attention \citep{arivazhagan2019monotonic,ma2020monotonic} or commit at meaning units \citep{zhang2020learning}; local agreement \citep{liu2020lowlatency,polak2022cuni,machacek2023turning} and attention \citep{papi2023alignatt} run offline models simultaneously, also on unbounded audio \citep{papi2024streamatt} and with LLM translators \citep{fuxa2026alignatt4llm}.
LLM-based systems translate streaming speech \citep{ouyang2025infinisst,guo2025streamuni,cheng2024towards}, and StreamSpeech, SeamlessStreaming and Hibiki(-Zero) also output speech \citep{zhang2024streamspeech,seamless2023seamless,labiausse2025highfidelity,labiausse2026simultaneous}.
All of them are evaluated per direction and without speaker attribution, the setting that SimulST toolkits and latency metrics assume \citep{ma2020simuleval,gaido2025simulstream,papi2022overgeneration,polak2025better}; several serve as our baselines (\cref{sec:exp}), and we reuse \laal{}, meaning units and Hibiki's contextual alignment (\cref{sec:bench,sec:model}).

\paragraph{Joint transcription, translation and speakers.}
Transcripts and translations can share one model \citep{chen2021direct,papi2023token}, and serialized output training transcribes overlapping speakers with one decoder, also streaming in time order (t-SOT) and with speaker attribution \citep{kanda2020serialized,kanda2022streaming,kanda2022streamingspeaker,kanda2020joint,kanda2021end}; \model reuses this time-ordered serialization.
Speaker embeddings let streaming ST detect speaker changes \citep{wang2025streaming}.
Of the closest systems (\cref{sec:intro}), DiariST builds on t-SOT and scores speaker-attributed BLEU on meetings translated from Chinese \citep{yang2024diarist}; STAC-ST obtains its two-speaker channel by merging the two sides of Fisher--CALLHOME calls \citep{zuluagagomez2023end}; JSTAR uses a fast--slow cascaded transducer encoder \citep{moritz2025transcribing}; and Seed LiveInterpret~2.0 \citep{cheng2025seed} is evaluated on RealSI clips cut from public videos \citep{cheng2024towards}.
\Cref{tab:positioning} contrasts their settings with ours; none of the four covers German or Korean.

\paragraph{Conversational translation systems and data.}
Conversational speech translation was studied as early as JANUS-II, which compared push-to-talk with cross-talk recording \citep{lavie1996translation}, but public test data still fall short (\cref{tab:corpora}): they are one-sided (MSLT; \citealp{federmann2016microsoft,federmann2017microsoft}), re-enacted (DRAL; \citealp{ward2023dialogs,ward2024dialogs,avila2023towards}), mediated by machine translation and recorded per participant (InCroMin; \citealp{cechovic2025corpus}), from one source language (Fisher--CALLHOME; \citealp{post2013improved}) or text-only (chat translation; \citealp{farajian2020findings,liang2021modeling}).
Code-switched speech alternates languages within, not between, speakers \citep{weller2022end,alastruey2023towards,zhou2025csdialogue}.
\bench therefore builds on TAXI \citep{bas2001taxi,reichel2016bas}, real German--English telephone dialogues with human translations of both sides (\cref{sec:bench}).

\paragraph{Diarization, turn-taking and multistream models.}
Neural diarization handles overlap, unknown speaker counts and streaming input \citep{fujita2019end,horiguchi2020end,xue2021online}; Sortformer orders speakers by arrival, also online \citep{park2025sortformer,medennikov2025streaming}; transducers emit speaker roles or turns \citep{elshafey2019joint,xia2022turn}.
Turn-taking models predict both speakers' voice activity \citep{ekstedt2022voice} or turn shifts from text \citep{ekstedt2020turngpt} but learn from monolingual dyads, even when multilingual \citep{inoue2024multilingual}.
Full-duplex dialogue models \citep{nguyen2023generative,defossez2024moshi,veluri2024beyond} and delayed or step-synchronous decoders \citep{zeghidour2025streaming,seide2024speech} run on a fixed clock.
Our diarization baselines are pyannote \citep{bredin2023pyannote} and Streaming Sortformer (\cref{sec:exp-aux}), and \model adopts arrival-order speakers and one clock-synchronous token stream (\cref{sec:model}).
\draftnote{planned: the baseline runs, and \model{} beyond its streaming ASR backbone}
